\subsection{The conductor graph and the horizontal spectral cover}

\begin{statement}{P4.55}{Proved}
In the remaining positive ruled case, the image $A$ of the
horizontal pole section is a smooth rational curve. The map
$H\to A$ has degree exactly two. The third generic
normalization preimage along $A$ is supplied by one rational
vertical conductor component, mapping isomorphically to $A$.
The other image component retains its unrestricted transverse
type. A separated self-identification point on either image
component is impossible.
\end{statement}
\begin{proof}
We write $T\to N\to S$ for the minimal normalization
resolution. By P4.54 and P4R.18, the pole section $H$ is
noncontracted, has genus $g\ge1$, and maps with degree two
or three to $A^\nu=\mathbb P^1$. The full generic transverse
germ along $A$ is $D_4$. The reduced image conductor is
$C=A\cup B$, with two rational normalized components and
one graph cycle.

The reduced normalization conductor $D$ has a tree as its
point-identification graph. Its positive-genus component
$H$ is smoothly embedded: the ruling descends to $N$, and
restriction of that ruling is an inverse to $H\to N$.
In particular no two points of $H$ are identified by the
normalization resolution. Every other normalized conductor
component is rational. These are the actual conductor
curves; vertical components and their identifications are
retained throughout the argument.

At any point of $C$, points of $N$ over it correspond to
local irreducible hypersurface branches of $S$. Two such
branches intersect along a curve, by the height bound for
two equations in a regular local ring of dimension three.
Their nonempty supports among the local analytic branches
of $C$ therefore pairwise intersect.

\paragraph{Separated self-identifications.}
Suppose $p\in B\setminus A$ has two analytic branches.
Let $D_B$ be the union of the conductor components mapping
to $B$, with $c_B$ normalized components and $s_B\ge1$
connected components. Its graph is a forest. Every component
maps finitely and surjectively to $B^\nu$, so the numbers
$r_+,r_-$ of its normalized points over the two preimages
of $p$ satisfy $r_+,r_-\ge c_B$.

There are $n$ normalization points of $S$ over $p$.
Their nonempty supports in $\{+,-\}$ pairwise intersect.
They have a common element: a singleton forces its element
into every support, and without a singleton all supports
are the entire two-element set. After interchanging signs,
$r_+\ge n$. The graph contribution at $p$ is consequently
\[
                  r_++r_--n\ge c_B.
\]
The whole forest has contribution $c_B-s_B<c_B$, a
contradiction. The same proof applies to a self-identification
of $A$ away from $B$, using all components mapping to $A$.
Thus the image cycle consists either of two distinct
intersections of $A$ and $B$, or of one three-branch point
belonging to both components.

\paragraph{The actual primary conductor and smoothness of $A$.}
Let $C_{A,\mathrm{sch}}$ be the primary component of the
actual image conductor supported on $A$. The primary
restriction maps of P4R.13, used in P4R.18, give
\[
                    g-1\le\chi(\mathcal O_{C_{A,\mathrm{sch}}}).
                                                            \tag{4.55a}
\]
The generic ordinary-triple conductor is the transverse
maximal ideal squared. Purity of the actual conductor
quotient identifies this primary ideal with $P_A^{(2)}$,
as in P4.31. The exact normalized conormal calculation of
P4.34 therefore gives
\[
 \chi(\mathcal O_{C_{A,\mathrm{sch}}})
       =5-\sum_p\bigl(\delta(A,p)+\tau_p-r_p\bigr).
                                                            \tag{4.55b}
\]
A planar point contributes exactly $5\delta(A,p)$, and a
nonplanar singular point with at most two branches contributes
at least seven. There cannot be three branches of $A$ at one
point: the identity
$\sum_p(r_p(C)-1)=2$ would then leave no identification for
the necessary connection with $B$; placing that connection
at the same point gives at least four branches and also
exceeds the identity.

Since $g\ge1$, (4.55a)--(4.55b) bound the total defect by
five. If $A$ is singular, it is planar everywhere, has total
delta invariant one, and $g=1$. It is then a locally complete
intersection curve in $Y$. Its conormal
$Q=P_A/P_A^2$ is locally free of rank two. Let
$j:\mathbb P^1\to A$ be its normalization and put $E_0=j^*Q$.
Adjunction gives
\[
             \omega_A=(\omega_Y|_A)\otimes(\det Q)^{-1}.
\]
The degree of $j^*\omega_Y$ is zero because $b_2(Y)=0$.
Also $p_a(A)=1$, so $\deg j^*\omega_A=0$. The degree of a
line bundle is preserved by finite birational normalization:
the normalization sequences for that line bundle and for
$\mathcal O_A$ have the same finite-length correction.
Consequently
\[
                              \deg E_0=0.                 \tag{4.55c}
\]

The actual reduced equation has generic order three along
$A$. Powers of a complete-intersection curve ideal are
unmixed, so $P_A^{(3)}=P_A^3$. Its third normal term is an
actual global section
\[
 q_3\in H^0\!\left(A,
          \operatorname{Sym}^3 Q\otimes\mathcal O_Y(S)|_A\right).
\]
The line twist has degree zero and becomes trivial on
$\mathbb P^1$. On the quotient projective bundle
$\mathbb P(E_0)$, with tautological class $\xi$ and fibre
class $F$, the pulled-back cubic has zero divisor of class
$3\xi$. Remove its common zero fibres, of total degree
$b\ge0$. The residual spectral curve $\Gamma$ has class
$3\xi-bF$. Its horizontal components are reduced, because
the generic $D_4$ cubic has three distinct roots.

We specify the spectral morphisms. At a generic point of
$H$, the first normal derivative of $T\to S\subset Y$
gives the direction of its smooth transverse branch and
hence a point of $\mathbb P(E_0)$. On the smooth compact
curve $H$, its two local coordinates extend after cancelling
their common power. This gives a holomorphic map over the
actual map $H\to\mathbb P^1$. Distinct generic preimages
on $H$ give distinct cubic roots. Hence its map onto the
spectral component $\Gamma_H$ is finite and birational,
and $H$ is that component's normalization. In degree three,
$\Gamma=\Gamma_H$. In degree two, precisely one remaining
rational vertical conductor component supplies the third
root, with degree one over $A$. Its direction curve
$\Gamma_R$ is a section and is isomorphic to $\mathbb P^1$.
Thus $\Gamma=\Gamma_H\cup\Gamma_R$. These exhaust the
three generic roots.

For a rank-two bundle of degree $e$ on $\mathbb P^1$,
\[
 \xi^2=e,\qquad \xi F=1,\qquad F^2=0,\qquad
 K_{\mathbb P(E_0)}=-2\xi+(e-2)F.
\]
Adjunction gives the exact arithmetic genus
\[
                    p_a(3\xi-bF)=3e-2-2b.               \tag{4.55d}
\]
With (4.55c) this is $-2-2b$. In degree three it is at
least $g\ge1$. In degree two the exact union sequence gives
\[
 p_a(\Gamma)=p_a(\Gamma_H)+p_a(\Gamma_R)
           +\operatorname{length}(\Gamma_H\cap\Gamma_R)-1
       \ge g-1\ge0.
\]
Both alternatives are impossible. Therefore $A$ is smooth.

\paragraph{The three-branch image point in degree three.}
Suppose now that $H\to A$ has degree three. Then $H$ is
the only conductor component mapping onto $A$. Since $A$
is smooth, a single three-branch image point consists of
one $A$ branch and two $B$ branches, denoted $+$ and $-$.

Let $c_B$ be the number of normalized conductor components
mapping to $B$. The entire conductor has $1+c_B$
normalized components and a tree graph, so its total
point-identification contribution is $c_B$. At the proposed
point write $r_A\ge1$, $r_+,r_-\ge c_B$ for the numbers
of conductor branches, and $n$ for the number of normalization
points. Each normalization point has at most one $A$
conductor branch, because $H$ is embedded.

A singleton support $\{A\}$ forces every point to carry
$A$. Then $n=r_A$, and the graph contribution is
$r_++r_-\ge2c_B>c_B$. A singleton $\{+\}$ forces
$r_+\ge n$, giving contribution at least
$r_A+r_-\ge1+c_B$; the same applies to $\{-\}$.
Thus every support has at least two elements. Consequently
$r_A+r_++r_-\ge2n$, and its graph contribution is at least
\[
       \frac{r_A+r_++r_-}{2}
                         \ge c_B+\frac{r_A}{2}>c_B.
\]
This is again impossible. Only two distinct $A$--$B$
junctions remain in degree three.

\paragraph{Both degree-three junctions are common cubic zeros.}
Let these junctions be $p,q$. If a local normalization
branch over $p$ has conductor only over $A$, pairwise
hypersurface intersections force every normalization point
over $p$ to carry $A$. Thus $n=r_A$, and the graph
contribution at $p$ is $r_B\ge c_B$, using the entire
budget. At $q$ some normalization point must carry both
$A$ and $B$: otherwise an $A$-only point and a $B$-only
point have disjoint supports. That point contributes
positively to the graph, contradicting the budget. There
is therefore no $A$-only local branch at either junction.

We retain all factors in the following local calculation.
Write the smooth curve as $A=\{x=y=0\}$. If its nonzero
cubic normal form has a simple root, use the normal blowup
chart $x=yu$. The implicit-function theorem on the strict
equation gives that root as $u=\alpha(y,z)$. The actual
function has the analytic factor
\[
                         f_1=x-y\alpha(y,z).
\]
This branch is smooth. Exact division writes $f=f_1f_2$,
and restriction of $f_2$ to this branch is $y^2$ times a
unit: the other cubic directions differ from the simple
root at the point. Its intersections with all other
branches are therefore supported only on $A$. It is the
forbidden $A$-only normalization branch.

Every nonzero binary cubic has a simple root unless it is
a cube. P4R.19 excludes a nonzero cube by its exact finite
weighted maps and compact intersection class. Hence the
whole cubic tensor vanishes at both $p$ and $q$.

\paragraph{The degree-three genus contradiction.}
For the now smooth curve $A=\mathbb P^1$, adjunction gives
$\deg N_{A/Y}=-2$, so $E=N_{A/Y}^*$ has degree two.
The actual cubic has at least two common zeros. Its residual
spectral divisor is $3\xi-bF$ with $b\ge2$, and is
irreducible with normalization $H$. Formula (4.55d) gives
\[
                    p_a(\Gamma)=4-2b\le0,
\]
contradicting $p_a(\Gamma)\ge g\ge1$. Degree three is
excluded. The degree-two alternative and its unique rational
third sheet are precisely the possibilities already identified
by the spectral morphism above.

The exact intersection polynomial and the graph inequalities
are recomputed in
\path{code/part4/d4_graph_spectral.py}; its saved output
reports PASS. The inequalities in this proof establish the
unrestricted graph counts.
\end{proof}
